\documentclass[11pt]{article}
\usepackage[a4paper,margin=1in]{geometry}
\usepackage{fontspec}
\usepackage[boldfont=false]{xeCJK}
\setCJKmainfont{FandolSong-Regular.otf}
\usepackage{amsmath}
\usepackage{amssymb}
\usepackage{array}
\usepackage{booktabs}
\usepackage{graphicx}
\usepackage{adjustbox}
\usepackage{enumitem}
\setlist[itemize]{topsep=2pt,partopsep=0pt,parsep=0pt,itemsep=2pt,leftmargin=1.4em}
\setlist[enumerate]{topsep=2pt,partopsep=0pt,parsep=0pt,itemsep=2pt,leftmargin=1.6em}
\usepackage{parskip}
\usepackage{fvextra}
\fvset{breaklines=true,breakanywhere=true,fontsize=\small}
\setlength{\parindent}{0pt}

\begin{document}

\begin{center}
  {\LARGE\bfseries Organisms and their environment}\\[0.35em]
  {\large Igcse Biology}
\end{center}
\vspace{0.6em}

\section*{Energy flow}

The \textbf{Sun} is the main source of \textbf{energy} 能量 for almost all life. Plants capture light energy by \textbf{photosynthesis} 光合作用 and store it as chemical energy in food. This energy then passes from organism to organism as they feed. At every step, some energy is lost (mostly as heat from \textbf{respiration} 呼吸作用) to the \textbf{environment} 环境. Energy flows \textbf{one way} — it is not recycled.

\section*{Food chains and food webs}

A \textbf{food chain} 食物链 shows the transfer of energy from one organism to the next, starting with a producer. Each arrow points in the direction the energy flows:

grass → rabbit → fox

A \textbf{food web} 食物网 is many food chains linked together.

\subsection*{Producers, consumers and decomposers}

\begin{center}
\adjustbox{max width=\linewidth}{%
\begin{tabular}{ll}
\toprule
\textbf{Type} & \textbf{Meaning} \\
\midrule
\textbf{producer} 生产者 & makes its own food (organic \textbf{nutrients} 营养物质), usually by photosynthesis \\
\textbf{consumer} 消费者 & gets its energy by feeding on other organisms \\
\textbf{herbivore} 食草动物 & a consumer that eats plants \\
\textbf{carnivore} 食肉动物 & a consumer that eats other animals \\
\textbf{decomposer} 分解者 & gets its energy from dead or waste material, causing \textbf{decomposition} 分解 \\
\bottomrule
\end{tabular}}
\end{center}

Consumers are named by their position: a \textbf{primary} consumer eats the producer; a \textbf{secondary} consumer eats the primary consumer; then come \textbf{tertiary} and \textbf{quaternary} consumers.

\subsection*{Trophic levels and pyramids}

A \textbf{trophic level} 营养级 is the position of an organism in a food chain (producers first, then primary consumers, and so on). You can show a food chain as a \textbf{pyramid} 金字塔:

\begin{itemize}
\item a \textbf{pyramid of numbers} counts the organisms at each level.

\item a \textbf{pyramid of biomass} 生物量 shows the total mass at each level — usually a better picture, because it does not depend on how big the organisms are.

\item \textbf{(Supplement)} a \textbf{pyramid of energy} shows the energy at each level — the most useful picture of all.

\end{itemize}

\subsection*{Energy loss along a chain (Supplement)}

Only about 10\% of the energy at one trophic level passes to the next. The rest is lost as heat (from respiration), in movement, and in waste. Because so much energy is lost, food chains usually have \textbf{fewer than five} trophic levels. It is more energy-efficient for people to eat \textbf{crop plants} 农作物 directly than to eat animals that were fed on those crops, because each extra level wastes energy.

\subsection*{Human impact}

Humans can damage food webs by:

\begin{itemize}
\item \textbf{overharvesting} 过度捕捞 — taking too many of one species (such as overfishing), so its numbers crash.

\item introducing a \textbf{foreign species} 外来物种 — a new species may have no predators and may crowd out native species.

\end{itemize}

\section*{Nutrient cycles}

Unlike energy, nutrients are \textbf{recycled} — used again and again.

\subsection*{The carbon cycle}

In the \textbf{carbon cycle} 碳循环:

\begin{itemize}
\item photosynthesis removes \textbf{carbon dioxide} 二氧化碳 from the air.

\item respiration and \textbf{combustion} 燃烧 (burning) return carbon dioxide to the air.

\item feeding passes carbon from one organism to the next.

\item decomposition returns carbon to the air and soil.

\item over millions of years, dead organisms can form \textbf{fossil fuels} 化石燃料 (coal and oil), which release carbon dioxide when they are burned.

\end{itemize}

\subsection*{The nitrogen cycle (Supplement)}

The \textbf{nitrogen cycle} 氮循环 recycles nitrogen for making \textbf{proteins} 蛋白质. \textbf{Microorganisms} 微生物 do most of the work:

\begin{itemize}
\item \textbf{decomposition}: decomposers break dead protein down into \textbf{ammonium ions} 铵离子.

\item \textbf{nitrification} 硝化作用: bacteria change ammonium ions into \textbf{nitrate ions} 硝酸根离子.

\item \textbf{nitrogen fixation} 固氮作用: lightning and some \textbf{bacteria} 细菌 turn nitrogen gas into compounds that plants can use.

\item plants absorb nitrate ions to make \textbf{amino acids} 氨基酸 and proteins; animals get them by feeding and \textbf{digestion} 消化.

\item \textbf{deamination} 脱氨基作用 (in the liver) and decomposition release nitrogen compounds again.

\item \textbf{denitrification} 反硝化作用: some bacteria turn nitrate ions back into nitrogen gas.

\end{itemize}

\section*{Populations}

\begin{itemize}
\item a \textbf{population} 种群 is a group of organisms of one \textbf{species} 物种, living in the same area at the same time.

\item a \textbf{community} 群落 is all the populations of different species in one place.

\item an \textbf{ecosystem} 生态系统 is the community together with its \textbf{environment}, all interacting.

\end{itemize}

\subsection*{Population growth}

The growth of a population depends on the food supply, \textbf{competition} 竞争, \textbf{predation} 捕食 (being eaten) and disease.

When a population grows in a place with limited \textbf{resources} 资源, its growth follows an S-shaped curve with four phases:

\begin{center}
\adjustbox{max width=\linewidth}{%
\begin{tabular}{ll}
\toprule
\textbf{Phase} & \textbf{What happens} \\
\midrule
\textbf{lag phase} 延迟期 & slow growth at first, while numbers are still small \\
\textbf{exponential phase} 对数期 & very fast growth, with plenty of food and space \\
\textbf{stationary phase} 稳定期 & growth stops; births ≈ deaths (a \textbf{limiting factor} 限制因素 such as food holds it back) \\
\textbf{death phase} 衰亡期 & numbers fall as food runs out, wastes build up, or disease spreads \\
\bottomrule
\end{tabular}}
\end{center}

\section*{Exam tips}

\begin{itemize}
\item Energy enters from the Sun, flows one way along the chain, and is lost as heat at each step. Nutrients are recycled.

\item Food chain starts with a producer; arrows show energy flow. Learn producer, consumer, herbivore, carnivore, decomposer.

\item Only \textasciitilde{}10\% of energy passes up each trophic level, so chains are short and eating plants is more efficient.

\item Carbon cycle: photosynthesis ↔ respiration and combustion; nitrogen cycle (Supplement): fixation, nitrification, denitrification by microorganisms.

\item Population: one species, one area, one time. Learn the four phases of the S-shaped growth curve.

\end{itemize}


\end{document}
